\chapter*{Tóm tắt}
\addcontentsline{toc}{chapter}{Tóm tắt}
\markboth{Tóm tắt}{Tóm tắt}

Báo cáo trình bày \textbf{Meeting AI Agent}, một hệ thống xử lý audio cuộc họp theo định hướng LLMOps. Hệ thống nhận file âm thanh, tạo transcript theo lượt nói, sinh tóm tắt, trích xuất action items có cấu trúc và cho phép người dùng hiệu chỉnh kết quả. Mục tiêu chính là chuyển nội dung cuộc họp từ dữ liệu không có cấu trúc thành thông tin có thể truy vết, đánh giá và tích hợp vào quy trình vận hành.

Kiến trúc hiện tại gồm Next.js frontend, FastAPI backend, Celery worker, Redis, PostgreSQL, Ollama/Qwen2.5-3B, Prometheus và Grafana. Audio được chuẩn hoá, đưa qua WhisperX để nhận dạng tiếng nói tự động (ASR), sau đó tổ chức thành transcript turns và truyền cho LLM để tạo summary hoặc action items dạng JSON. Trước khi lưu kết quả, hệ thống áp dụng guardrails như kiểm tra schema, lọc prompt injection, kiểm tra ngày hạn, phát hiện dấu hiệu hallucination và che một số thông tin nhạy cảm bằng regex-based PII masking.

LLMOps được thể hiện qua cơ chế feedback, lưu vết dữ liệu, benchmark, monitoring và quy trình chuẩn bị cho fine-tuning. Khi người dùng sửa task description, assignee, due date, đánh dấu false positive hoặc bổ sung task bị bỏ sót, các hiệu chỉnh này được lưu trong PostgreSQL và có thể export phục vụ đánh giá hoặc huấn luyện sau này. Fine-tuning hiện được xem là luồng đã chuẩn bị về kỹ thuật, chưa phải bằng chứng chính cho chất lượng baseline.

Đánh giá baseline được thực hiện trên 100 mẫu từ tập tham chiếu synthetic \texttt{data/eval/gold\_synthetic\_205.jsonl}, gồm transcript cuộc họp tổng hợp và action items kỳ vọng. Với Qwen2.5-3B ở chế độ few-shot, hệ thống đạt precision 0.8604, recall 0.6665, F1 0.6886, assignee accuracy 0.5232, hallucination rate 0.0\% và schema failure rate 0.0\%. Kết quả cho thấy pipeline ổn định về định dạng và kiểm soát hallucination trên benchmark tổng hợp, trong khi recall và assignee accuracy vẫn là các hướng cần cải thiện.
